\begin{proof}[Proof of \cref{obj:lem:exact-reflection-budget}]
We first prove the component estimate at \(s=1\), treating the one-dimensional and two-dimensional cases in that order, and then prove the estimate for \(0<s<1\).

\begin{enumerate}
\item \textit{Normalization and inverse truncations.}
For \(p\in\{1,2\}\), write
\[
d\mu_p(y)=2^{-p}\,dy,
\qquad
(\mathscr D_pG)(y)=G(b_p(y)),
\qquad
\widehat\zeta(a)=
\int_{\mathbb T_2^p}\zeta(y)e^{-\mathrm i\pi a^{\mathsf T}y}\,d\mu_p(y),
\]
where \(G\in L^2(\mathbb R^p;\mathbb C)\), \(\zeta\in L^2(\mathbb T_2^p;\mathbb C)\), and \(a\in\mathbb Z^p\). These operators act on equivalence classes. Indeed, for \(\epsilon\in\{0,1\}^p\), define the reflection branch
\[
\bigl(\mathscr R_\epsilon(x)\bigr)_j
=
\begin{cases}
x_j,&\epsilon_j=0,\\
2-x_j,&\epsilon_j=1,
\end{cases}
\qquad x\in[0,1]^p.
\]
Each branch preserves Lebesgue measure between the corresponding boxes, and
\(b_p(\mathscr R_\epsilon(x))=x\) almost everywhere. Consequently, for every nonnegative measurable \(\varphi\) on the cube,
\[
\int_{\mathbb T_2^p}\varphi(b_p(y))\,d\mu_p(y)
=
2^{-p}\sum_{\epsilon\in\{0,1\}^p}
\int_{[0,1]^p}\varphi(x)\,dx
=
\int_{[0,1]^p}\varphi(x)\,dx.
\]
In particular, folding preserves null sets in the required direction and
\[
\|\mathscr D_pG\|_{L^2(\mu_p)}^2
=
\int_{[0,1]^p}|G(u)|^2\,du
\leq
\|G\|_{L^2(\mathbb R^p)}^2.
\]
The normalized characters form a complete orthonormal Fourier basis, so
\[
\sum_{a\in\mathbb Z^p}|\widehat\zeta(a)|^2
=
\|\zeta\|_{L^2(\mu_p)}^2.
\]

Use the angular Fourier transform
\[
(\mathcal F_pG)(\omega)
=
(2\pi)^{-p/2}
\int_{\mathbb R^p}G(u)e^{-\mathrm i\omega^{\mathsf T}u}\,du
\]
on \(L^1\cap L^2\), with its unitary extension to \(L^2\).
For a measurable profile \(\psi\in L^2(\mathbb R^p;\mathbb C)\) and an integer
\(n_{\mathrm{cut}}\geq0\), define
\[
\psi_{n_{\mathrm{cut}}}(\omega)
=
\mathbf1_{\{\|\omega\|_2\leq n_{\mathrm{cut}}\}}\psi(\omega),
\qquad
G_{\psi,n_{\mathrm{cut}}}(u)
=
(2\pi)^{-p/2}
\int_{\mathbb R^p}
e^{\mathrm i\omega^{\mathsf T}u}
\psi_{n_{\mathrm{cut}}}(\omega)\,d\omega.
\]
Every frequency moment of the truncated profile is integrable. Explicitly, for every integer \(\ell_{\mathrm{mom}}\geq0\), Cauchy--Schwarz gives
\[
\int_{\mathbb R^p}
\|\omega\|_2^{\ell_{\mathrm{mom}}}
|\psi_{n_{\mathrm{cut}}}(\omega)|\,d\omega
\leq
\left(
\int_{\{\|\omega\|_2\leq n_{\mathrm{cut}}\}}
\|\omega\|_2^{2\ell_{\mathrm{mom}}}\,d\omega
\right)^{1/2}
\|\psi\|_2
<\infty.
\]
At \(n_{\mathrm{cut}}=0\), the cutoff is zero almost everywhere because \(p>0\).
Plancherel and spatial negation give
\[
\|G_{\psi,n_{\mathrm{cut}}}\|_2^2
=
\|\psi_{n_{\mathrm{cut}}}\|_2^2,
\qquad
\|G_{\psi,n_{\mathrm{cut}}}-G_{\psi,n'_{\mathrm{cut}}}\|_2^2
=
\|\psi_{n_{\mathrm{cut}}}-\psi_{n'_{\mathrm{cut}}}\|_2^2.
\]
Dominated convergence makes these inverse truncations Cauchy in \(L^2\).

To identify their limit explicitly, let \(\mathcal F_{\mathrm{cyc},p}\) denote the unitary \(L^2\) extension of
\[
(\mathcal F_{\mathrm{cyc},p}\psi)(\eta)
=
\int_{\mathbb R^p}\psi(\omega)
e^{-2\pi\mathrm i\eta^{\mathsf T}\omega}\,d\omega,
\qquad \psi\in L^1\cap L^2,\quad \eta\in\mathbb R^p,
\]
and define, as an \(L^2\) class,
\[
(\mathcal I_p\psi)(u)
=
(2\pi)^{-p/2}
(\mathcal F_{\mathrm{cyc},p}\psi)\!\left(-\frac{u}{2\pi}\right).
\]
The squared prefactor \((2\pi)^{-p}\) cancels the dilation Jacobian \((2\pi)^p\). Thus
\[
\|\mathcal I_p\psi\|_2^2=\|\psi\|_2^2,
\qquad
\|\mathcal I_p\psi-\mathcal I_p\chi\|_2^2=\|\psi-\chi\|_2^2,
\qquad \psi,\chi\in L^2(\mathbb R^p;\mathbb C).
\]
The defining integral identifies
\(G_{\psi,n_{\mathrm{cut}}}=\mathcal I_p\psi_{n_{\mathrm{cut}}}\) almost everywhere, whence
\[
\|G_{\psi,n_{\mathrm{cut}}}-\mathcal I_p\psi\|_2^2
=
\|\psi_{n_{\mathrm{cut}}}-\psi\|_2^2
\longrightarrow0.
\]

For an admissible extension \(G\in L^1\cap L^2\), its original spatial class is recovered by this construction. Fubini, applied to the integrable product of the truncated profile and an \(L^1\cap L^2\) test function \(\varphi\), gives
\[
\int_{\mathbb R^p}\overline{\varphi(u)}
G_{\psi,n_{\mathrm{cut}}}(u)\,du
=
\int_{\mathbb R^p}
\overline{(\mathcal F_p\varphi)(\omega)}
\psi_{n_{\mathrm{cut}}}(\omega)\,d\omega.
\]
Taking \(\psi=\mathcal F_pG\) and \(\varphi=G\), and expanding the squared distance using this pairing and the two Plancherel identities, yields
\[
\|G-G_{\mathcal F_pG,n_{\mathrm{cut}}}\|_2^2
=
\|\mathcal F_pG-(\mathcal F_pG)_{n_{\mathrm{cut}}}\|_2^2
\longrightarrow0.
\]
Uniqueness of the \(L^2\) limit therefore gives
\[
\mathcal I_p(\mathcal F_pG)=G
\quad\text{almost everywhere}.
\]
% lean: CausalSmith.Experimentation.BivariateSobolevDesignCapacity.reflectionInverseL2Representative_fourier_ae_original

\item \textit{The one-dimensional endpoint.}
Differentiating the inverse truncation under the integral is justified by the frequency moments established above. For \(1\leq j\leq p\),
\[
\partial_jG_{\psi,n_{\mathrm{cut}}}(u)
=
(2\pi)^{-p/2}
\int_{\mathbb R^p}
\mathrm i\omega_j e^{\mathrm i\omega^{\mathsf T}u}
\psi_{n_{\mathrm{cut}}}(\omega)\,d\omega.
\]
These derivatives are continuous and square integrable. Applying the inverse-integral energy identity to the profile
\(\mathrm i\omega_j\psi_{n_{\mathrm{cut}}}\) gives
\[
\int_{\mathbb R^p}
|\partial_jG_{\psi,n_{\mathrm{cut}}}(u)|^2\,du
=
\int_{\mathbb R^p}
\omega_j^2|\psi_{n_{\mathrm{cut}}}(\omega)|^2\,d\omega.
\]
Summing the coordinate identities proves the exact first-order identity
\[
\int_{\mathbb R^p}
\left(
|G_{\psi,n_{\mathrm{cut}}}(u)|^2+
\frac1p\sum_{j=1}^p
|\partial_jG_{\psi,n_{\mathrm{cut}}}(u)|^2
\right)\,du
=
\int_{\mathbb R^p}
\left(1+\frac{\|\omega\|_2^2}{p}\right)
|\psi_{n_{\mathrm{cut}}}(\omega)|^2\,d\omega.
\]

Now take \(p=1\), identifying a one-dimensional integer frequency with \(a\in\mathbb Z\). Define the signed reflected derivative on the period-two circle by
\[
V_{\psi,n_{\mathrm{cut}},1}(y)
=
\begin{cases}
G'_{\psi,n_{\mathrm{cut}}}(y),&0<y<1,\\
-G'_{\psi,n_{\mathrm{cut}}}(2-y),&1<y<2,\\
0,&y\in\{0,1\},
\end{cases}
\]
using representatives in \([0,2)\). For a continuously differentiable periodic test function \(\varphi\), integration by parts on the two intervals has total boundary term
\[
\bigl[G_{\psi,n_{\mathrm{cut}}}(y)\varphi(y)\bigr]_0^1
+
\bigl[G_{\psi,n_{\mathrm{cut}}}(2-y)\varphi(y)\bigr]_1^2
=
G_{\psi,n_{\mathrm{cut}}}(0)\bigl(\varphi(2)-\varphi(0)\bigr)
=
0.
\]
It follows that
\[
\int_{\mathbb T_2^1}
(\mathscr D_1G_{\psi,n_{\mathrm{cut}}})(y)\varphi'(y)\,d\mu_1(y)
=
-\int_{\mathbb T_2^1}
V_{\psi,n_{\mathrm{cut}},1}(y)\varphi(y)\,d\mu_1(y).
\]
Taking \(\varphi(y)=e^{-\mathrm i\pi ay}\) gives the exact multiplier identity
\[
\widehat V_{\psi,n_{\mathrm{cut}},1}(a)
=
\mathrm i\pi a\,
\widehat{\mathscr D_1G_{\psi,n_{\mathrm{cut}}}}(a).
\]
Parseval for the reflected function and its signed derivative, followed by folding, therefore gives
\[
\begin{aligned}
\sum_{a\in\mathbb Z}(1+\pi^2a^2)
\left|\widehat{\mathscr D_1G_{\psi,n_{\mathrm{cut}}}}(a)\right|^2
&=
\int_0^1
\left(
|G_{\psi,n_{\mathrm{cut}}}(u)|^2+
|G'_{\psi,n_{\mathrm{cut}}}(u)|^2
\right)\,du\\
&\leq
\int_{\mathbb R}
\left(
|G_{\psi,n_{\mathrm{cut}}}(u)|^2+
|G'_{\psi,n_{\mathrm{cut}}}(u)|^2
\right)\,du\\
&=
\int_{\mathbb R}(1+\omega^2)
|\psi_{n_{\mathrm{cut}}}(\omega)|^2\,d\omega\\
&\leq
\int_{\mathbb R}(1+\omega^2)|\psi(\omega)|^2\,d\omega.
\end{aligned}
\]

For an admissible extension \(G\), take \(\psi=\mathcal F_1G\).
The \(L^2\) contraction of \(\mathscr D_1\) and the spatial convergence above imply
\[
\mathscr D_1G_{\mathcal F_1G,n_{\mathrm{cut}}}
\longrightarrow\mathscr D_1G
\quad\text{in }L^2(\mu_1).
\]
Every Fourier coefficient converges, since its absolute difference is bounded by the \(L^2\) distance. For every finite frequency set \(\mathcal A_{\mathrm{fin}}\subset\mathbb Z\), passing to the limit in its finite weighted sum gives
\[
\sum_{a\in\mathcal A_{\mathrm{fin}}}(1+\pi^2a^2)
|\widehat{\mathscr D_1G}(a)|^2
\leq
\int_{\mathbb R}(1+\omega^2)|\mathcal F_1G(\omega)|^2\,d\omega.
\]
Taking the supremum over finite sets proves the same inequality for the full sum. Folding transports \(G=g\) almost everywhere on the cube to
\(\mathscr D_1G=h\) almost everywhere on the torus. Infimizing over admissible extensions gives
\[
\sum_{a\in\mathbb Z}(1+\pi^2a^2)|\widehat h(a)|^2
\leq N_{1,1}(g)^2.
\]
Extensions with infinite first-order energy satisfy the intermediate inequality in the extended sense as well.
% lean: CausalSmith.Experimentation.BivariateSobolevDesignCapacity.componentFourierBudget_one_firstOrder_le

\item \textit{The two-dimensional endpoint.}
For \(p=2\), define, for \(j\in\{1,2\}\),
\[
\sigma_j(y)=
\begin{cases}
1,&0<y_j<1,\\
-1,&1<y_j<2,\\
0,&y_j\in\{0,1\},
\end{cases}
\qquad
V_{\psi,n_{\mathrm{cut}},j}(y)
=
\sigma_j(y)
(\partial_jG_{\psi,n_{\mathrm{cut}}})(b_2(y)).
\]
Apply the one-dimensional integration-by-parts computation on each coordinate slice and integrate in the other coordinate. The boundary terms cancel on each slice, giving
\[
\int_{\mathbb T_2^2}
(\mathscr D_2G_{\psi,n_{\mathrm{cut}}})(y)
\partial_j\varphi(y)\,d\mu_2(y)
=
-\int_{\mathbb T_2^2}
V_{\psi,n_{\mathrm{cut}},j}(y)\varphi(y)\,d\mu_2(y).
\]
For \(a=(a_1,a_2)\in\mathbb Z^2\), the character test yields
\[
\widehat V_{\psi,n_{\mathrm{cut}},j}(a)
=
\mathrm i\pi a_j
\widehat{\mathscr D_2G_{\psi,n_{\mathrm{cut}}}}(a).
\]
Joint Parseval, obtained from the product Fourier basis, combines the zeroth-order energy and the two derivative energies with factor \(1/2\):
\[
\begin{aligned}
&\sum_{a\in\mathbb Z^2}
\left(1+\frac{\pi^2(a_1^2+a_2^2)}2\right)
\left|\widehat{\mathscr D_2G_{\psi,n_{\mathrm{cut}}}}(a)\right|^2\\
&\quad=
\int_{[0,1]^2}
\left(
|G_{\psi,n_{\mathrm{cut}}}(u)|^2+
\frac{
|\partial_1G_{\psi,n_{\mathrm{cut}}}(u)|^2+
|\partial_2G_{\psi,n_{\mathrm{cut}}}(u)|^2}{2}
\right)\,du\\
&\quad\leq
\int_{\mathbb R^2}
\left(
|G_{\psi,n_{\mathrm{cut}}}(u)|^2+
\frac{
|\partial_1G_{\psi,n_{\mathrm{cut}}}(u)|^2+
|\partial_2G_{\psi,n_{\mathrm{cut}}}(u)|^2}{2}
\right)\,du\\
&\quad=
\int_{\mathbb R^2}
\left(1+\frac{\|\omega\|_2^2}{2}\right)
|\psi_{n_{\mathrm{cut}}}(\omega)|^2\,d\omega\\
&\quad\leq
\int_{\mathbb R^2}
\left(1+\frac{\|\omega\|_2^2}{2}\right)|\psi(\omega)|^2\,d\omega.
\end{aligned}
\]
For \(\psi=\mathcal F_2G\), spatial \(L^2\) convergence and the reflection contraction again give convergence of every reflected coefficient. Passing first to each finite frequency sum and then to their supremum proves
\[
\sum_{a\in\mathbb Z^2}
\left(1+\frac{\pi^2\|a\|_2^2}{2}\right)
|\widehat{\mathscr D_2G}(a)|^2
\leq
\int_{\mathbb R^2}
\left(1+\frac{\|\omega\|_2^2}{2}\right)
|\mathcal F_2G(\omega)|^2\,d\omega.
\]
The equality \(\mathscr D_2G=h\) almost everywhere for every admissible extension, followed by the extension infimum, yields
\[
\sum_{a\in\mathbb Z^2}
\left(1+\frac{\pi^2\|a\|_2^2}{2}\right)|\widehat h(a)|^2
\leq N_{2,1}(g)^2.
\]
% lean: CausalSmith.Experimentation.BivariateSobolevDesignCapacity.componentFourierBudget_two_firstOrder_le

\item \textit{The common operator for fractional smoothness.}
Assume now that \(0<s<1\). On the complete frequency \(L^2\) space, define
\[
\mathscr T_p\psi=\mathscr D_p(\mathcal I_p\psi),
\qquad p\in\{1,2\}.
\]
This is a common linear operator for the two endpoints. For \(\vartheta\in[0,1]\), define the extended gauges
\[
\|\psi\|_{\mathrm{freq},p,\vartheta}^2
=
\int_{\mathbb R^p}
\left(1+\frac{\|\omega\|_2^2}{p}\right)^\vartheta
|\psi(\omega)|^2\,d\omega,
\]
\[
\|\zeta\|_{\mathrm{per},p,\vartheta}^2
=
\sum_{a\in\mathbb Z^p}
\left(1+\frac{\pi^2\|a\|_2^2}{p}\right)^\vartheta
|\widehat\zeta(a)|^2,
\]
for \(\psi\in L^2(\mathbb R^p;\mathbb C)\) and
\(\zeta\in L^2(\mathbb T_2^p;\mathbb C)\); a divergent expression has value \(+\infty\).

The inverse isometry and folding prove
\[
\|\mathscr T_p\psi\|_{\mathrm{per},p,0}
\leq
\|\psi\|_{\mathrm{freq},p,0}.
\]
For every frequency \(L^2\) class, the inverse truncations converge to
\(\mathcal I_p\psi\), so
\[
\mathscr D_pG_{\psi,n_{\mathrm{cut}}}
\longrightarrow\mathscr T_p\psi
\quad\text{in }L^2(\mu_p).
\]
Apply the smooth spectral estimate in the appropriate dimension and pass to finite sums as above. This proves the second endpoint on the same operator:
\[
\|\mathscr T_p\psi\|_{\mathrm{per},p,1}
\leq
\|\psi\|_{\mathrm{freq},p,1}.
\]
The inequality includes profiles with infinite first-order gauge.
% lean: CausalSmith.Experimentation.BivariateSobolevDesignCapacity.reflectionProfileOperator_firstOrder_norm_le

\item \textit{The normalized interpolation energy.}
For two extended gauges \(\mathsf n_0,\mathsf n_1\) on a common additive space, define, for \(\rho>0\),
\[
\mathfrak K^2(\rho;\mathsf n_0,\mathsf n_1;\zeta)
=
\inf_{\zeta=\zeta_0+\zeta_1}
\left\{
\mathsf n_0(\zeta_0)^2+
\rho^2\mathsf n_1(\zeta_1)^2
\right\},
\]
and, for \(0<s<1\),
\[
\mathfrak J_s(\mathsf n_0,\mathsf n_1;\zeta)
=
\frac{2\sin(\pi s)}{\pi}
\int_0^\infty
\rho^{-1-2s}
\mathfrak K^2(\rho;\mathsf n_0,\mathsf n_1;\zeta)\,d\rho.
\]
Every source decomposition maps to a target decomposition under \(\mathscr T_p\). The two endpoint contractions therefore imply
\[
\begin{aligned}
&\mathfrak K^2\!\left(
\rho;
\|\cdot\|_{\mathrm{per},p,0},
\|\cdot\|_{\mathrm{per},p,1};
\mathscr T_p\psi
\right)\\
&\qquad\leq
\mathfrak K^2\!\left(
\rho;
\|\cdot\|_{\mathrm{freq},p,0},
\|\cdot\|_{\mathrm{freq},p,1};
\psi
\right).
\end{aligned}
\]
Integration against the positive normalized measure gives the same contraction for \(\mathfrak J_s\).

We evaluate the source functional exactly. Set
\[
\mathsf w_{\mathrm{freq},p}(\omega)
=
1+\frac{\|\omega\|_2^2}{p},
\qquad
\psi_{0,\rho}(\omega)
=
\frac{\rho^2\mathsf w_{\mathrm{freq},p}(\omega)}
{1+\rho^2\mathsf w_{\mathrm{freq},p}(\omega)}\psi(\omega),
\]
\[
\psi_{1,\rho}(\omega)
=
\frac{1}{1+\rho^2\mathsf w_{\mathrm{freq},p}(\omega)}\psi(\omega).
\]
These are an admissible \(L^2\) decomposition of \(\psi\). Their cost satisfies
\[
|\psi_{0,\rho}(\omega)|^2+
\rho^2\mathsf w_{\mathrm{freq},p}(\omega)
|\psi_{1,\rho}(\omega)|^2
=
\frac{\rho^2\mathsf w_{\mathrm{freq},p}(\omega)}
{1+\rho^2\mathsf w_{\mathrm{freq},p}(\omega)}
|\psi(\omega)|^2
\leq|\psi(\omega)|^2.
\]
For every other decomposition
\(\psi=\chi_0+\chi_1\), completing the square gives
\[
\begin{aligned}
|\chi_0|^2+\rho^2\mathsf w_{\mathrm{freq},p}|\chi_1|^2
&=
\frac{\rho^2\mathsf w_{\mathrm{freq},p}}
{1+\rho^2\mathsf w_{\mathrm{freq},p}}|\psi|^2\\
&\quad+
(1+\rho^2\mathsf w_{\mathrm{freq},p})
\left|
\chi_1-\frac{\psi}{1+\rho^2\mathsf w_{\mathrm{freq},p}}
\right|^2.
\end{aligned}
\]
Thus the displayed decomposition attains the infimum, including when the competing decomposition has infinite cost, and
\[
\mathfrak K^2\!\left(
\rho;
\|\cdot\|_{\mathrm{freq},p,0},
\|\cdot\|_{\mathrm{freq},p,1};
\psi
\right)
=
\int_{\mathbb R^p}
\frac{\rho^2\mathsf w_{\mathrm{freq},p}(\omega)}
{1+\rho^2\mathsf w_{\mathrm{freq},p}(\omega)}
|\psi(\omega)|^2\,d\omega.
\]
For every \(\beta>0\), substitution
\(\tau=\rho\sqrt{\beta}\) gives
\[
\frac{2\sin(\pi s)}{\pi}
\int_0^\infty
\rho^{-1-2s}\frac{\rho^2\beta}{1+\rho^2\beta}\,d\rho
=
\beta^s\frac{2\sin(\pi s)}{\pi}
\int_0^\infty\frac{\tau^{1-2s}}{1+\tau^2}\,d\tau
=
\beta^s,
\]
because the last integral is \(\pi/(2\sin(\pi s))\).
Tonelli consequently proves
\[
\mathfrak J_s\!\left(
\|\cdot\|_{\mathrm{freq},p,0},
\|\cdot\|_{\mathrm{freq},p,1};
\psi
\right)
=
\|\psi\|_{\mathrm{freq},p,s}^2.
\]
Combining this identity with the operator contraction yields
\[
\mathfrak J_s\!\left(
\|\cdot\|_{\mathrm{per},p,0},
\|\cdot\|_{\mathrm{per},p,1};
\mathscr T_p\psi
\right)
\leq
\|\psi\|_{\mathrm{freq},p,s}^2.
\]
% lean: CausalSmith.Experimentation.BivariateSobolevDesignCapacity.reflectionProfileOperator_kNormSq_le_energy

\item \textit{Transport through Fourier coefficients.}
Define the coefficient map and its sequence gauges by
\[
\mathscr C_p\zeta=(\widehat\zeta(a))_{a\in\mathbb Z^p},
\qquad
\mathsf w_{\mathrm{per},p}(a)
=
1+\frac{\pi^2\|a\|_2^2}{p},
\]
\[
\|\gamma\|_{\mathrm{seq},p,0}^2
=
\sum_{a\in\mathbb Z^p}|\gamma(a)|^2,
\qquad
\|\gamma\|_{\mathrm{seq},p,1}^2
=
\sum_{a\in\mathbb Z^p}
\mathsf w_{\mathrm{per},p}(a)|\gamma(a)|^2,
\qquad \gamma\in\ell^2(\mathbb Z^p;\mathbb C).
\]
Parseval and the first-order gauge definition give
\[
\|\mathscr C_p\zeta\|_{\mathrm{seq},p,0}
=
\|\zeta\|_{\mathrm{per},p,0},
\qquad
\|\mathscr C_p\zeta\|_{\mathrm{seq},p,1}
=
\|\zeta\|_{\mathrm{per},p,1}.
\]
Mapping each torus decomposition through \(\mathscr C_p\) therefore gives
\[
\mathfrak J_s\!\left(
\|\cdot\|_{\mathrm{seq},p,0},
\|\cdot\|_{\mathrm{seq},p,1};
\mathscr C_p\zeta
\right)
\leq
\mathfrak J_s\!\left(
\|\cdot\|_{\mathrm{per},p,0},
\|\cdot\|_{\mathrm{per},p,1};
\zeta
\right).
\]
The same pointwise minimizing decomposition and scalar integral from the previous step, now with counting measure and weight
\(\mathsf w_{\mathrm{per},p}\), show that
\[
\mathfrak J_s\!\left(
\|\cdot\|_{\mathrm{seq},p,0},
\|\cdot\|_{\mathrm{seq},p,1};
\mathscr C_p\zeta
\right)
=
\sum_{a\in\mathbb Z^p}
\mathsf w_{\mathrm{per},p}(a)^s|\widehat\zeta(a)|^2.
\]
All these weights are finite and at least one, so the minimizing decompositions belong to the ambient \(\ell^2\) space. Applying the two inequalities to \(\zeta=\mathscr T_p\psi\) yields
\[
\sum_{a\in\mathbb Z^p}
\left(1+\frac{\pi^2\|a\|_2^2}{p}\right)^s
|\widehat{\mathscr T_p\psi}(a)|^2
\leq
\int_{\mathbb R^p}
\left(1+\frac{\|\omega\|_2^2}{p}\right)^s
|\psi(\omega)|^2\,d\omega.
\]

For an admissible extension \(G\), the inverse identification in the first step gives
\[
\mathscr T_p(\mathcal F_pG)=\mathscr D_pG=h
\quad\text{almost everywhere}.
\]
Hence
\[
\sum_{a\in\mathbb Z^p}
\left(1+\frac{\pi^2\|a\|_2^2}{p}\right)^s
|\widehat h(a)|^2
\leq
\int_{\mathbb R^p}
\left(1+\frac{\|\omega\|_2^2}{p}\right)^s
|\mathcal F_pG(\omega)|^2\,d\omega.
\]
Infimizing over all admissible extensions proves the component assertion for \(0<s<1\). Together with the two endpoint cases, this establishes it throughout \(0<s\leq1\).
% lean: CausalSmith.Experimentation.BivariateSobolevDesignCapacity.componentFourierBudget_le_sobolevNormSq

\item \textit{Coefficient additivity and the order-two support bound.}
Fix \(d\geq2\) and \(m\in\mathcal H_{d,s}\). By \cref{obj:def:sobolev-class}, \(m\) is centered and square integrable and satisfies \cref{obj:ass:order-two,obj:ass:pooled-budget}. Use the canonical components defined in \cref{obj:synth_1}. Define their ambient reflections by
\[
h_j^{\mathrm{amb}}(y)
=
g_j(m)((b_d(y))_j),
\qquad
h_{j\ell}^{\mathrm{amb}}(y)
=
g_{j\ell}(m)((b_d(y))_j,(b_d(y))_\ell),
\qquad j<\ell,
\]
and their coefficients by
\[
\widehat h_j^{\mathrm{amb}}(k)
=
2^{-d}\int_{[0,2)^d}
h_j^{\mathrm{amb}}(y)e^{-\mathrm i\pi k^{\mathsf T}y}\,dy,
\]
\[
\widehat h_{j\ell}^{\mathrm{amb}}(k)
=
2^{-d}\int_{[0,2)^d}
h_{j\ell}^{\mathrm{amb}}(y)e^{-\mathrm i\pi k^{\mathsf T}y}\,dy,
\qquad k\in\mathbb Z^d.
\]

The requisite component \(L^2\) regularity follows from conditional Cauchy--Schwarz. In the integration notation of \cref{obj:synth_1},
\[
\int_0^1|g_j(m)(x_j)|^2\,dx_j
\leq
\int_{[0,1]^d}|m(x)|^2\,dx,
\]
and
\[
\int_{[0,1]^2}
\left|
\int_{[0,1]^{d-2}}
m(x_j,x_\ell,u_{-\{j,\ell\}})\,du_{-\{j,\ell\}}
\right|^2\,dx_j\,dx_\ell
\leq
\int_{[0,1]^d}|m(x)|^2\,dx.
\]
Subtracting the two lifted main effects therefore gives an \(L^2\) pair effect. For \(d=2\), the inner zero-dimensional integral is evaluation, and the second inequality is equality. Folding preserves this regularity. The almost-everywhere decomposition supplied by \cref{obj:ass:order-two} thus permits integration term by term:
\[
\widehat F(k)
=
\sum_{j=1}^d\widehat h_j^{\mathrm{amb}}(k)
+
\sum_{j<\ell}\widehat h_{j\ell}^{\mathrm{amb}}(k).
\]

To establish coordinate locality, fix a coordinate \(\jmath\) with
\(k_\jmath\neq0\), and define a torus translation by
\[
(\tau_{k,\jmath})_\iota
=
\begin{cases}
-1/k_\jmath\pmod{2},&\iota=\jmath,\\
0\pmod{2},&\iota\neq\jmath,
\end{cases}
\qquad 1\leq\iota\leq d.
\]
If a reflected component is independent of coordinate \(\jmath\), it is invariant under this translation, whereas its coefficient character acquires the factor
\[
e^{-\mathrm i\pi k^{\mathsf T}\tau_{k,\jmath}}=-1.
\]
Translation invariance of Haar measure then makes the component coefficient equal to its own negative, hence zero. In particular,
\[
\widehat h_j^{\mathrm{amb}}(k)=0
\quad\text{if some }k_\jmath\neq0\text{ with }\jmath\neq j,
\]
\[
\widehat h_{j\ell}^{\mathrm{amb}}(k)=0
\quad\text{if some }k_\jmath\neq0\text{ with }\jmath\notin\{j,\ell\}.
\]
If \(|\operatorname{supp}(k)|>2\), choose three distinct active coordinates
\(\jmath_1,\jmath_2,\jmath_3\). Every main component excludes at least one of them, and every pair component also excludes at least one. All terms in the coefficient decomposition vanish, proving
\[
\widehat F(k)=0
\qquad\text{when }|\operatorname{supp}(k)|>2.
\]
% lean: CausalSmith.Experimentation.BivariateSobolevDesignCapacity.Fhat_zero_of_orderTwo_from_components

\item \textit{Centering and unique component coefficients.}
Fubini applied to the canonical marginal formulas gives
\[
\int_0^1g_j(m)(\upsilon)\,d\upsilon
=
\int_{[0,1]^d}m(x)\,dx
=
0.
\]
Integrating the raw pair marginal in its second coordinate recovers the first main marginal. Consequently, for almost every \(x_j\),
\[
\int_0^1g_{j\ell}(m)(x_j,\upsilon)\,d\upsilon
=
g_j(m)(x_j)-g_j(m)(x_j)
-\int_0^1g_\ell(m)(\upsilon)\,d\upsilon
=
0.
\]
Interchanging the two marginal coordinates gives, for almost every \(x_\ell\),
\[
\int_0^1g_{j\ell}(m)(\upsilon,x_\ell)\,d\upsilon
=
g_\ell(m)(x_\ell)
-\int_0^1g_j(m)(\upsilon)\,d\upsilon
-g_\ell(m)(x_\ell)
=
0.
\]
These identities also apply at \(d=2\), with the evaluation convention for the raw pair marginal.

Folding preserves the coordinate integrals. Thus the main reflected coefficient at zero vanishes. For a pair effect, if \(k_\ell=0\), integrating first in coordinate \(\ell\) gives zero by the first pair-centering identity; if \(k_j=0\), use the other identity. Combining these facts with coordinate locality yields
\[
\widehat h_j^{\mathrm{amb}}(k)\neq0
\ \Longrightarrow\
\operatorname{supp}(k)=\{j\},
\]
\[
\widehat h_{j\ell}^{\mathrm{amb}}(k)\neq0
\ \Longrightarrow\
\operatorname{supp}(k)=\{j,\ell\}.
\]
The supports of distinct canonical components are therefore disjoint. Substituting these support statements into coefficient additivity proves
\[
\widehat F(k)
=
\begin{cases}
\widehat h_j^{\mathrm{amb}}(k),
&\operatorname{supp}(k)=\{j\},\\
\widehat h_{j\ell}^{\mathrm{amb}}(k),
&\operatorname{supp}(k)=\{j,\ell\},\quad j<\ell.
\end{cases}
\]
For the pair case, ordering the two supporting coordinates increasingly selects the unique pair in the decomposition.
% lean: CausalSmith.Experimentation.BivariateSobolevDesignCapacity.Fhat_orderTwo_pair_component

\item \textit{Pooling and local coefficient transport.}
For the auxiliary sums over all ambient frequencies, extend the stated complexity by
\[
t(0)=0,
\qquad
t(k)=\frac{\|k\|_2^2}{|\operatorname{supp}(k)|}
\quad(k\neq0).
\]
At a singleton-support frequency, the unique main coefficient just identified supplies the full coefficient; at a pair-support frequency, the unique ordered pair supplies it. At every other frequency the indicator below is zero. Hence, pointwise in \(k\),
\[
\begin{aligned}
&\mathbf1_{\{1\leq|\operatorname{supp}(k)|\leq2\}}
(1+\pi^2t(k))^s|\widehat F(k)|^2\\
&\quad\leq
\sum_{j=1}^d
(1+\pi^2t(k))^s|\widehat h_j^{\mathrm{amb}}(k)|^2
+
\sum_{j<\ell}
(1+\pi^2t(k))^s|\widehat h_{j\ell}^{\mathrm{amb}}(k)|^2.
\end{aligned}
\]
Summing and interchanging the nonnegative frequency sums with the finite component sums gives
\[
\begin{aligned}
&\sum_{\substack{k\in\mathbb Z^d\\1\leq|\operatorname{supp}(k)|\leq2}}
(1+\pi^2t(k))^s|\widehat F(k)|^2\\
&\quad\leq
\sum_{j=1}^d\sum_{k\in\mathbb Z^d}
(1+\pi^2t(k))^s|\widehat h_j^{\mathrm{amb}}(k)|^2\\
&\qquad+
\sum_{j<\ell}\sum_{k\in\mathbb Z^d}
(1+\pi^2t(k))^s|\widehat h_{j\ell}^{\mathrm{amb}}(k)|^2.
\end{aligned}
\]

Define the local reflections and coefficients by
\[
h_j^{\mathrm{loc}}(y)
=
g_j(m)(b_1(y)),
\qquad
\widehat h_j^{\mathrm{loc}}(\alpha)
=
\frac12\int_{[0,2)}
h_j^{\mathrm{loc}}(y)e^{-\mathrm i\pi\alpha y}\,dy,
\qquad \alpha\in\mathbb Z,
\]
\[
h_{j\ell}^{\mathrm{loc}}(y)
=
g_{j\ell}(m)((b_2(y))_1,(b_2(y))_2),
\qquad
\widehat h_{j\ell}^{\mathrm{loc}}(a)
=
\frac14\int_{[0,2)^2}
h_{j\ell}^{\mathrm{loc}}(y)e^{-\mathrm i\pi a^{\mathsf T}y}\,dy,
\qquad a\in\mathbb Z^2.
\]
For \(\alpha\in\mathbb Z\) and \(a\in\mathbb Z^2\), define the ambient frequency embeddings
\[
(\iota_j(\alpha))_\jmath
=
\begin{cases}
\alpha,&\jmath=j,\\
0,&\jmath\neq j,
\end{cases}
\qquad
(\iota_{j\ell}(a))_\jmath
=
\begin{cases}
a_1,&\jmath=j,\\
a_2,&\jmath=\ell,\\
0,&\jmath\notin\{j,\ell\}.
\end{cases}
\]
Coordinate projection preserves product Haar probability, and folding commutes with selecting coordinates. Integrating the unused coordinates therefore gives
\[
\widehat h_j^{\mathrm{amb}}(\iota_j(\alpha))
=
\widehat h_j^{\mathrm{loc}}(\alpha),
\qquad
\widehat h_{j\ell}^{\mathrm{amb}}(\iota_{j\ell}(a))
=
\widehat h_{j\ell}^{\mathrm{loc}}(a).
\]
At the frequencies where the corresponding ambient coefficients can be nonzero,
\[
t(\iota_j(\alpha))=\alpha^2
\quad(\alpha\neq0),
\qquad
t(\iota_{j\ell}(a))=\frac{a_1^2+a_2^2}{2}
\quad(a_1a_2\neq0).
\]
Indeed, the ambient squared frequency length is the local squared length, and the support cardinality is respectively one or two.

The ambient support statements restrict each sum to these embedded frequencies. Restriction to the supporting coordinates is injective, so enlarging to all local frequencies gives
\[
\sum_{k\in\mathbb Z^d}
(1+\pi^2t(k))^s|\widehat h_j^{\mathrm{amb}}(k)|^2
\leq
\sum_{\alpha\in\mathbb Z}
(1+\pi^2\alpha^2)^s|\widehat h_j^{\mathrm{loc}}(\alpha)|^2,
\]
\[
\sum_{k\in\mathbb Z^d}
(1+\pi^2t(k))^s|\widehat h_{j\ell}^{\mathrm{amb}}(k)|^2
\leq
\sum_{a\in\mathbb Z^2}
\left(1+\frac{\pi^2\|a\|_2^2}{2}\right)^s
|\widehat h_{j\ell}^{\mathrm{loc}}(a)|^2.
\]
Thus the ambient bound is bounded in turn by the sum of the local component Fourier budgets.
% lean: CausalSmith.Experimentation.BivariateSobolevDesignCapacity.reflectedBudget_le_local_component_budgets

\item \textit{The pooled budget and cube energy.}
The local canonical components are Borel and square integrable by the marginal argument above. Their restriction norms are finite because the finite, nonnegative sum in \cref{obj:ass:pooled-budget} is bounded by one. Applying the component contraction to each main effect and each pair effect yields
\[
\begin{aligned}
&\sum_{\substack{k\in\mathbb Z^d\\1\leq|\operatorname{supp}(k)|\leq2}}
(1+\pi^2t(k))^s|\widehat F(k)|^2\\
&\quad\leq
\sum_{j=1}^d\sum_{\alpha\in\mathbb Z}
(1+\pi^2\alpha^2)^s|\widehat h_j^{\mathrm{loc}}(\alpha)|^2\\
&\qquad+
\sum_{j<\ell}\sum_{a\in\mathbb Z^2}
\left(1+\frac{\pi^2\|a\|_2^2}{2}\right)^s
|\widehat h_{j\ell}^{\mathrm{loc}}(a)|^2\\
&\quad\leq
\sum_{j=1}^dN_{1,s}(g_j(m))^2
+
\sum_{j<\ell}N_{2,s}(g_{j\ell}(m))^2
\leq1,
\end{aligned}
\]
where the final inequality is precisely \cref{obj:ass:pooled-budget}.
% lean: CausalSmith.Experimentation.BivariateSobolevDesignCapacity.reflectedBudget_le_one_of_sobolevClass

Folding sends Haar probability to cube Lebesgue probability in dimension \(d\), by the same coordinatewise reflection calculation used above. Centering therefore gives
\[
\widehat F(0)
=
\int_{\mathbb T_2^d}F(y)\,d\mu_d(y)
=
\int_{[0,1]^d}m(x)\,dx
=
0,
\qquad d\mu_d(y)=2^{-d}\,dy.
\]
Since \(d\geq2\), Parseval applies in positive dimension. Also, \(s>0\) and \(t(k)\geq0\) give
\[
(1+\pi^2t(k))^s\geq1.
\]
Using the zero coefficient and the already established vanishing for support size greater than two, we conclude
\[
\begin{aligned}
\int_{[0,1]^d}|m(x)|^2\,dx
&=
\int_{\mathbb T_2^d}|F(y)|^2\,d\mu_d(y)\\
&=
\sum_{k\in\mathbb Z^d}|\widehat F(k)|^2\\
&=
\sum_{\substack{k\in\mathbb Z^d\\1\leq|\operatorname{supp}(k)|\leq2}}
|\widehat F(k)|^2\\
&\leq
\sum_{\substack{k\in\mathbb Z^d\\1\leq|\operatorname{supp}(k)|\leq2}}
(1+\pi^2t(k))^s|\widehat F(k)|^2
\leq1.
\end{aligned}
\]
This establishes all the asserted conclusions.
% lean: CausalSmith.Experimentation.BivariateSobolevDesignCapacity.cube_second_moment_le_of_reflectedBudget
\end{enumerate}
\end{proof}
